\documentclass[12pt]{article}
\title{Exp6: Polarization and Crystal Display\\
Final Report}
\author{Group1}
\linespread{1.1}
\setlength{\parskip}{5mm}

\begin{document}
\maketitle

\section*{Theoretical analysis}

If we put a $ \lambda /4$   waveplate between parallel linear polarizers which directions are $ x $ polarized and $ y $ polarized, and the incident laser beam  passing the first polarizer denotes $ \mathbf{J}=[0,1]^T $ by Jones calculus. Assume $ \theta $ is the angle between fast axis(or slow axis$^\dag$ ) and laboratory coordinate $ x $-axis, and $ \beta $ is the rotating angle of the analyzer (the second polarizer).

After passing $ \lambda /4$   waveplate, the Jones vector becomes


\begin{eqnarry*}
$\mathbf{J}'$ &=& $\mathbf{M}$_{\lambda /4}\mathbf{J}=\mathbf{R}(\theta)\left[\begin{array}{cc}e^{-i\Gamma /2}&0\\0&e^{i\Gamma /2}\end{array}\right]\mathbf{R}(-\theta)\mathbf{J}\\
&=&\left[\begin{array}{cc}e^{-i\Gamma /2}\cos^2 \theta+e^{i\Gamma /2}\sin^2 \theta&-i\sin (\Gamma /2)\sin 2\theta&-i\sin (\Gamma/2)\sin  2\theta&e^{-i\Gamma /2}\sin^2 \theta+e^{i\Gamma /2}\cos^2 \theta\end{array}\right]\mathbf{J}
\end{eqnarray*}

With $ \Gamma = \pi /4 $ and $ \mathbf{J}=[0,1]^T $, we have

\begin{align}
$\mathbf{J}'$=\left[\begin{array}{c}-i\sin (\Gamma /2)\sin 2\theta&e^{-i\Gamma /2}\sin^2 \theta+e^{i\Gamma /2}\cos^2 \theta\end{array}\right]=\frac{1}{\sqrt{2}}\left[\begin{array}{c}-i\sin 2\theta &1+i\cos 2\theta\end{array}\right]
\end{align}

Passing the analyzer, the final Jones vector is

\begin{align}
$\mathbf{J}_o$=$\mathbf{R}$(-\beta)\left[\begin{array}{cc}1&0&0&0\end{array}\right]\mathbf{R}(\beta)\mathbf{J}=\frac{1}{\sqrt{2}}\left[ \begin{array}{c}
\sin\beta\cos\beta+i(\sin 2\beta\cos 2 \theta /2-\cos^2\beta\sin 2\theta)&\sin^2\beta+i(\sin^2\beta\cos 2\theta -\sin 2\beta\sin 2 \theta /2)\end{array}\right]
\end{align}

The optical intensity is proportional to $ \mathbf{J}^\dag\mathbf{J} $. Obviously the intensity is a function of $ \theta $ and $ \beta $.

\begin{align}
2I= 2$ \mathbf{J}_o^\dag\mathbf{J}_o$ = \sin^2\beta (1+\cos^2 2\theta)+\cos^2\beta \sin^2 2 \theta +\sin 2\theta\cos 2\theta\sin 2 \beta
\end{align}

With trivial caculations,
\begin{eqnarray*}
2I&=&\sin^2 \beta +(\cos\beta\sin 2 \theta +\sin \beta\cos 2\theta)^2=\sin^2\beta+\sin^2(\beta +2\theta)\\
&=&1-\frac{1}{2}(\cos 2\beta +\cos (2\beta +4\theta))=1-\frac{1}{2}\left\lbrace \cos 2\beta [1+\cos 4\beta]+\sin 2\beta\sin 4\beta\right\rbrace \\
&=&1-\frac{1}{2}A\cos (2\beta +\alpha)
\end{eqnarry*}


where

$$
A=\sqrt{(1+\cos 4\theta)^2+\sin^2 4\theta}=\sqrt{2+ 2\cos 4\theta}=2\cos 2\theta
$$

$$
\alpha =-\arctan \left(\frac{\sin 4\theta}{1+\cos 4\theta}\right)
$$

The output intensity is a sinusoidal function

$$
2I=1-\cos 2\theta\cos (2\beta +\alpha)
$$

The maximum corresponds with the long axis of polarization ellipse,
which means we could measure the elliptic polarized light produced by $ \lambda $/4 waveplate.
$ \alpha $ is the azumith of  polarization ellipse which could be calculated from $ \theta $ or varing $ \beta $.

 
\section*{Consider special case(keep $ \lambda $/4 waveplate unmoved)}



\subsection*{$ \theta $ = (3n+1)$ \pi $/4}

$$
2I=\sin^2\beta +\cos^2\beta =1
$$

Since the output beam of $ \theta $ = $ \pi $/4 waveplate  are cicular polarized,
the intensity drop to const with ratio 1/2 regardless of the angle $ \beta $.

\subsection*{$ \theta $ = n$ \pi $/2}

$$
2I=2\sin^2 \beta
$$

Since the input beam's polarizaion direction is parallel to slow axis or fast axis,
the output polariization suffers zerp phase different between two  eigenbasis.
The intensity varies just as Mulus law.
\end{document}